
\section{「正确」态射定义的动机}

这里给出两个动机，说明态射应当有怎样的行为。第一个是代数的，第二个是几何的。

\textbf{7.1.1.}\ \emph{代数动机。}\quad 我们将会希望：诸\trAffineScheme{}的%
\trMorphism{} $\Spec A\to\Spec B$ 恰好就是诸\trRingMorphism{} $B\to A$。我们已经%
见过，\trRingMorphism{} $B\to A$ 会朝相反的方向诱导诸\trTopologicalSpace{}之间的%
映射（\cref{ex:ring-maps-continuous-maps}）；主要的新配料，则是要掺进这份混合物%
里的函数\trStructureSheaf。于是，概形之间的\trMorphism{}就应当是这样一个东西——%
「在\trAffine\trOpenSet{}的层面上，看上去便是这副模样。」

\textbf{7.1.2.}\ \emph{几何动机。}\quad 受\trManifold{}理论（\S3.1.1）的驱动——%
\trManifold{}与\trScheme{}一样，都是\trRingedSpace——我们希望概形之间的%
\trMorphism{}至少得是\trRingedSpace{}之间的\trMorphism；至于这类态射为何物，我们%
现在就来为它寻觅动机。（我们将在下一节把它形式化。）注意：若 $\pi\colon X\to Y$ 是%
诸\trDifferentiable\trManifold{}之间的一个映射，则 $Y$ 上的一个光滑函数拉回到
$X$ 上，得到的仍是 $X$ 上的一个光滑函数。更精确地说，给定一个\trOpenSubset{}
$\mathsf{U}\subset Y$，就有一个自然的映射
$\Gamma(\mathsf{U},\OO_Y)\to\Gamma(\pi^{-1}(\mathsf{U}),\OO_X)$。这条映射与限制%
良好相容（把一个函数限制到更小的\trOpenSet{}上再拉回，与先拉回再限制，结果相同），
所以事实上我们有 $Y$ 上诸\trSheaf{}的一个映射 $\OO_Y\to\pi_*\OO_X$。同样地，概形%
之间的一个\trMorphism{} $\pi\colon X\to Y$ 也应当诱导一个映射
$\OO_Y\to\pi_*\OO_X$。但事实上，在诸\trDifferentiable\trManifold{}的\trCategory{}%
中，一个\trContinuousMap{} $\pi\colon X\to Y$ 是诸\trDifferentiable\trManifold{}%
之间的一个映射，当且仅当 $Y$ 上的光滑函数拉回到 $X$ 上仍是光滑函数（也就是说，从
$Y$ 上的光滑函数到 $X$ 上的\emph{函数}的那个\trPullback{}映射，实际上正落在%
\emph{光滑函数}这个子集之中；也就是说，\trContinuousMap{} $\pi$ 诱导了光滑函数的%
一个\trPullback，而它可以解读为一个映射 $\OO_Y\to\pi_*\OO_X$）。于是，这个%
\trSheaf{}的映射便\emph{刻画}了\trDifferentiable{}范畴中的态射。所以，我们本可以%
拿它当\trDifferentiable{}范畴中\trMorphism{}的\emph{定义}来用（见%
\cref{ex:pullback-detects-smoothness}）。

但我们要怎样把这套道理用到概形的\trCategory{}上去呢？在诸\trDifferentiable%
\trManifold{}的\trCategory{}中，一个\trContinuousMap{} $\pi\colon X\to Y$ 会%
\emph{诱导}函数之\trSheaf{}的一个\trPullback，我们可以问：它何时会诱导\emph{光滑}%
函数的一个\trPullback。然而，函数在概形上要更为古怪，单凭诸\trTopologicalSpace{}%
之间的映射，我们复原不出那个\trPullback{}映射。正确的补丁是把这一点硬性焊进%
\trMorphism{}的定义之中：也就是说，取一个\trContinuousMap{} $\pi\colon X\to Y$，
连同一条\trPullback{}映射 $\pi^{\sharp}\colon\OO_Y\to\pi_*\OO_X$。由此便得到赋环%
空间的\emph{\trCategory}的定义。

人们或许指望把概形之间的\trMorphism{}定义为\trRingedSpace{}之间的\trMorphism。
但这并不全对，因为那样一来，动机~7.1.1就不满足了：如我们所愿，对每个\trMorphism{}
$A\to B$ 都有\trMorphism{} $\Spec B\to\Spec A$，可是还会冒出额外的、并非这样产生
的\trRingedSpace{}的\trMorphism{} $\Spec B\to\Spec A$（见\mbox{练习~7.2.G}）。把
定义改成「局部看起来具有这种形状的\trRingedSpace{}之间的\trMorphism」也能行得通；
但为了避免梳理诸般细节的麻烦，我们改走另一条路——使用\trLocallyRingedSpace{}，
这相当于不仅要求函数拉回，还要求\emph{函数的值}拉回。不过我们将验证：我们最终的%
定义实际上与它等价（\mbox{练习~7.3.D}）。

我们先来正式地定义\trRingedSpace{}的\trMorphism。
